\section{The Effective Gradient Ratio Diagnostic}
\label{sec:egr}

Theorems~\ref{thm:hessian} and~\ref{thm:twoscale} together identify a
training pathology that is invisible to standard loss curves: as
cross-entropy decays toward zero on both training and validation
sets, the spectral module is silently gradient-starved. A practitioner
running joint training sees the loss go down and the validation
accuracy go up, with no in-training signal that the spectral module
has failed to specialize. The implication for held-out test
performance is only revealed once test data is examined---often after
many epochs of wasted compute, or after deployment.

This section introduces the \emph{Effective Gradient Ratio} (EGR), a
real-time diagnostic that detects gradient starvation as it occurs.
EGR is a single scalar, computable inside any training loop at
negligible cost, that converts the finite-time starvation prediction
of Theorem~\ref{thm:twoscale} into a concrete observable.

\subsection{Definition and basic properties}

Recall from Definition~\ref{def:egr} that at training step $t$,
\[
    \EGR(t) = \frac{\norm{\nabla_\thetap \loss(\thetap_t, \phip_t)}}
                   {\norm{\nabla_\phip \loss(\thetap_t, \phip_t)}}.
\]
The numerator and denominator are computable directly from gradients
accumulated during the standard backward pass. EGR is unitless and
invariant to the choice of learning rate.

\paragraph{Initial value: a heuristic.}
A useful rule of thumb for the starting point of the diagnostic is
\begin{equation}
\label{eq:egr_init}
    \EGR(0) \;\sim\; \sqrt{\Cf / \Cg}
    \qquad \text{(heuristic, not a theorem).}
\end{equation}
The reasoning is a parameter-counting one: at variance-scaled
initialization each module's gradient has entries of comparable
typical magnitude, so the squared gradient norms scale with the
number of parameters carrying them,
$\norm{\nabla_\phip \loss}^2 \sim \Cg$ against
$\norm{\nabla_\thetap \loss}^2 \sim \Cf$, and the square root comes
from passing to norms. We state~\eqref{eq:egr_init} as a heuristic
rather than a proposition for three reasons: it exchanges the
expectation of a ratio for the ratio of expectations without a
concentration argument; the per-entry-magnitude premise is
parametrization-dependent and would change under $\mu$P, which we
cite elsewhere~\cite{yang2021tensor}; and the spectral module's
gradient additionally passes through the bottleneck, which the count
ignores. Its role in this paper is limited to setting the scale for
the \emph{relative} threshold of Definition~\ref{def:alert}, which
compares $\EGR(t)$ against the measured $\EGR(0)$ of the run at hand
and so does not depend on~\eqref{eq:egr_init} being sharp. The
corresponding numerical value for our motivating architecture follows
from the capacity ratio reported in Section~\ref{sec:real}.

\subsection{Decay of the spectral gradient}

\begin{proposition}[Spectral gradient decay]
\label{prop:egr_monotonic}
Under the hypotheses of Theorem~\ref{thm:twoscale}
(Assumptions~\ref{ass:linear}, \ref{ass:inputlip},
and~\ref{ass:residual}),
\[
    \norm{\nabla_\thetap \loss(t)}
    \;\leq\; \frac{B\,C}{1 + \mu t},
\]
where $B$ is the operator-norm cap of Lemma~\ref{lem:opcap} and
$C, \mu$ are the constants of Assumption~\ref{ass:residual}.
\end{proposition}

\begin{proof}
By the chain rule~\eqref{eq:chain},
$\norm{\nabla_\thetap \loss} \leq \norm{J_\thetap}_{\mathrm{op}}
\norm{r} \leq B \norm{r}$, and Assumption~\ref{ass:residual} bounds
$\norm{r(t)}$.
\end{proof}

\begin{remark}[From gradient decay to the EGR ratio]
\label{rem:egr_ratio_caveat}
Proposition~\ref{prop:egr_monotonic} bounds the \emph{numerator} of
the EGR. The denominator $\norm{\nabla_\phip \loss} =
\norm{J_\phip^\top r}$ is also proportional to the residual, so both
quantities vanish together as the loss is fitted, and the behavior of
their ratio is not determined by residual decay alone: it depends on
the relative conditioning of the two Jacobians along the trajectory.
We therefore treat the EGR trajectory as an empirical object,
characterized in Section~\ref{sec:egr-empirical}, and use
Proposition~\ref{prop:egr_monotonic} only for the absolute claim that
the spectral module's learning signal decays polynomially while the
loss is being fitted.
\end{remark}

\subsection{Threshold criterion}

Proposition~\ref{prop:egr_monotonic} suggests a practical criterion
for detecting the onset of gradient starvation.

\begin{definition}[Gradient starvation alert]
\label{def:alert}
We say the model is in a \emph{gradient starvation regime} at step $t$
if
\[
    \EGR(t) \;\;<\;\; \tau \cdot \EGR(0),
\]
where $\tau \in (0, 1)$ is a user-supplied tolerance. Empirically, we
recommend $\tau \in [0.05, 0.1]$, corresponding to a 10--20$\times$
drop in EGR relative to initialization.
\end{definition}

\begin{remark}
The threshold $\tau$ is the only hyperparameter of the diagnostic.
The default $\tau = 0.1$ corresponds to a one-decade decay and is the
value we use throughout Section~\ref{sec:experiments}. We provide a
sensitivity analysis to $\tau$ in the supplement.
\end{remark}

\paragraph{Practical use.}
A practitioner logs $\EGR(t)$ at each step. When the alert fires,
two options are available:
\begin{enumerate}
    \item \textbf{Freeze the spectral module.} This is the prescription
          most directly motivated by Theorem~\ref{thm:twoscale}:
          removing $\thetap$ from the optimization eliminates the
          timescale gap. We discuss this in Section~\ref{sec:discussion}.
    \item \textbf{Inject a fixed residual baseline.} A softer
          intervention: replace $Z = f_\thetap(X)$ with
          $Z = f_\thetap(X) + h(X)$ for a fixed feature extractor
          $h$ (e.g.\ PCA). The fixed baseline ensures the spatial
          model receives stable input even if $f_\thetap$ continues
          to drift, breaking the shortcut without freezing.
\end{enumerate}
We evaluate both interventions in Experiment~1.3.

\subsection{Relationship to existing diagnostics}

EGR is conceptually related to several existing notions but is
distinct in its scope and intent.

\paragraph{GradNorm~\cite{chen2018gradnorm}.}
GradNorm dynamically reweights losses in multi-task learning to
equalize gradient magnitudes across tasks. EGR measures rather than
intervenes, and applies to single-task compositional architectures
rather than multi-task losses.

\paragraph{Modality-specific learning rates (MSLR).}
In multimodal learning, MSLR adjusts learning rates per modality. EGR
diagnoses the underlying pathology that MSLR addresses but in a
unimodal, two-stage architecture where the ``modalities'' are
the spectral and spatial subspaces of a single input.

\paragraph{Loss curve monitoring.}
The conventional practice of monitoring training and validation loss
fails to detect gradient starvation because both can decay even as
$\thetap$ stagnates. EGR is sensitive to the structural failure
even when no global loss curve abnormality is present.

\subsection{Empirical validation: EGR as a diagnostic}
\label{sec:egr-empirical}

We validate EGR as a diagnostic of training pathology in two stages.
First, we establish that EGR statistics correlate with generalisation
across our full sweep. Second, we test whether the signal is genuine
or merely a re-encoding of model capacity.

\paragraph{Pooled correlation (capacity confounded).}
Across $N=24$ runs spanning four widths $\{16,64,256,1024\}$ with six
seeds each, EGR depth correlates strongly with final test accuracy:
Pearson $r = -0.723$ ($p = 6.6\times 10^{-5}$), Spearman $\rho = -0.808$
($p = 1.8\times 10^{-6}$), and with the train--test gap at $r = +0.741$
($p = 3.5\times 10^{-5}$). However, EGR depth is itself strongly
correlated with $\log\mathrm{width}$ ($r = +0.903$): roughly $77\%$ of
the pooled correlation with test accuracy is attributable to capacity
rather than to EGR depth per se (partial $r$ controlling for
$\log\mathrm{width}$ collapses to $-0.169$, $p = 0.44$). The pooled
correlation alone cannot establish a capacity-independent diagnostic.

\paragraph{Fixed-capacity test.}
To isolate the capacity-independent signal, we ran a second experiment
holding CNN width fixed at $D = 256$ and varying data noise
$\{0.05,0.10,0.15,0.20\}$, learning rate
$\{3{\times}10^{-4},10^{-3},3{\times}10^{-3}\}$, weight decay
$\{0,10^{-3}\}$, and six seeds, for a total of $N = 144$ runs. At fixed
capacity, $\mathrm{EGR}_{\min}$ is the relevant statistic (with width
fixed, $\mathrm{egr\_depth}$ reduces to a constant offset of
$\mathrm{egr\_min}$ in expectation). It exhibits a within-capacity
correlation with generalisation that is statistically robust but
practically moderate: Pearson $r(\mathrm{EGR}_{\min}, \mathrm{acc}) =
+0.344$ ($p = 2.5\times 10^{-5}$).

The relationship survives every adversarial control we applied.
$\mathrm{EGR}_{\min}$ is essentially orthogonal to the swept
hyperparameters --- marginal correlations with noise, learning rate,
and weight decay are $+0.029$, $-0.073$, and $+0.007$ respectively,
with the full hyperparameter set explaining only $R^{2} = 0.015$ of
$\mathrm{EGR}_{\min}$ variance. The partial correlation
$r(\mathrm{EGR}_{\min}, \mathrm{acc} \mid \mathrm{noise}, \mathrm{lr},
\mathrm{wd}) = +0.395$ ($p = 9.5\times 10^{-7}$) is \emph{larger} than
the pooled $r = +0.334$, which is the opposite signature of a hidden
hyperparameter confound. The within-bin analysis is the strongest
evidence: across the 24 $(\mathrm{noise}, \mathrm{lr}, \mathrm{wd})$
cells in which only the seed varies, 23 of 24 yield strongly positive
$r(\mathrm{EGR}_{\min}, \mathrm{acc})$ (mean $+0.584$, Fisher-$z$
aggregated $t = 14.65$, $p < 10^{-13}$).

Controlling for final training loss --- itself a non-trivial predictor
with $r = +0.47$ for accuracy and $r = -0.55$ for the overfit gap ---
leaves the $\mathrm{EGR}_{\min}$ partial correlation essentially
unchanged at $+0.335$ ($p = 4.1\times 10^{-5}$). EGR therefore carries
information that train-loss monitoring does not.

\paragraph{Practical interpretation.}
Honesty about effect size is required. Treated as a binary alarm,
$\mathrm{EGR}_{\min}$ achieves ROC-AUC $= 0.612$ for flagging
below-median test accuracy, and $\mathrm{EGR}_{\mathrm{depth}}$
achieves only $0.538$ for high overfit gap. The best operating point
($\mathrm{EGR}_{\min} < 0.10$) catches $90\%$ of bad runs but also
flags $57\%$ of healthy runs --- no threshold delivers a usable
trade-off as a standalone indicator.

\paragraph{Verdict.}
At fixed capacity, EGR is a real, capacity-independent signal
correlated with generalisation. The within-capacity effect size is
small enough that EGR is best deployed alongside train-loss monitoring
as a complementary diagnostic, not as a standalone run-quality
indicator. For research and ablation work, EGR provides genuine
insight into the gradient starvation mechanism beyond what loss
monitoring reveals; for production monitoring of single training
runs, plain loss curves remain the primary signal and EGR a useful
adjunct.

\subsection{Implementation}

The EGR diagnostic is implemented as a PyTorch training callback. It
hooks into the standard training loop between the backward pass and
the optimizer step:
\begin{verbatim}
loss.backward()
egr_logger.log_step(global_step)   # <-- the only added line
optimizer.step()
optimizer.zero_grad()
\end{verbatim}
The callback identifies the spectral and spatial parameter groups via
\verb|model.spectral.parameters()| and \verb|model.spatial.parameters()|
(or any user-supplied naming convention) and records the gradient
norms and EGR at each step. We release the implementation as the
\verb|egr| package, available at the project repository.
\TODO{Phase 4, W12: prepare PyPI package skeleton. Add unit tests
for gradient norm computation and threshold detection.}
